\chapter{Building, configuring, releasing and deploying}
\label{ch:build_release_deployment}

This chapter describes the project's current build and deployment workflow.
The developer setup script runs the build, release, and deployment steps in
sequence. They remain separate operations: the build creates staged
applications, the release packages those files, and deployment installs a
release with configuration for a selected environment.

\section{Building and releasing}
\label{sec:build_release_deployment:build_release}

The build script compiles the C++ components and, unless disabled, the Java
applications. It runs the unit and integration tests unless the caller asks to
skip tests. The built applications are staged under \texttt{installed/};
instrumented builds and some target platforms use a separate staging
directory.

The release script packages the staged applications with the configuration
templates, database scripts, environment profiles, and deployment and
developer-environment scripts. The version comes from the project build
configuration, and the Git revision is included in the release name by
default. A manifest in the archive identifies its contents. The release is
therefore a specific package of built files and deployment material, rather
than a build directory copied directly to a host.

\section{TOML configuration}
\label{sec:build_release_deployment:toml}

The project uses TOML in two related ways. Component configuration templates
contain the settings read by individual services. Environment profiles supply
the values used to prepare those templates for a particular installation.
This separates the settings a component reads from the host names, addresses,
paths, and other values that vary between installations.

During deployment, the selected profile is flattened into placeholder names.
For example, a value in a section named \texttt{arbiter\_primary} with the key
\texttt{peer\_host} supplies the
\texttt{\$\{arbiter\_primary\_peer\_host\}} placeholder in a component
template. Deployment fails and identifies the missing value if a template
contains an unresolved placeholder. It generates self-signed TLS certificates
unless \texttt{--skip-certs} is given; \texttt{--force-certs} regenerates
certificates that already exist. Deployment can create the database and export
the credentials used by the authentication service; options allow those
database steps to be skipped when the database is managed separately.

The C++ component loaders read the TOML keys they use while loading a
configuration. They do not currently reject every unrecognised key. The
proposed check for unknown keys is documented separately and is not part of
the current behaviour.

\section{Environment profiles}
\label{sec:build_release_deployment:profiles}

The profiles under \texttt{environments/} describe development, test,
pre-production, and production deployments. The development profile is the
complete working example for a single-machine sandbox. The other committed
profiles describe full high-availability deployments on dedicated hosts, but
are examples rather than ready-to-use configurations: they contain placeholder
hostnames and other values that must be supplied for a real environment.
Deployment refuses a profile when a required template value is missing.

The profile defines deployment paths, database and TLS settings, component
instances, host and port values, and the startup order used by the developer
environment script. For example, \texttt{devenv.py} starts components in the
order declared in \texttt{[startup\_order]}, and stops them in reverse order.
The \texttt{dev.toml} profile is also the default when no environment is
specified for that script.

\section{Deployment and starting services}
\label{sec:build_release_deployment:deploy}

Deployment unpacks the release into its installation directory, expands the
component configuration templates from the selected profile, prepares TLS
certificates when requested, and runs the database setup and credential
export steps unless those are skipped. The deployment script reports missing
template values rather than writing configuration that still contains an
unresolved placeholder.

For a developer sandbox, \texttt{devenv.py} reads the environment profile,
starts the configured component processes in startup order, and records their
logs and process identifiers under paths from that profile. It can also stop,
restart, and report the status of those processes. Production deployment uses
the deployment script and its profile; the developer environment script is
not the production process manager.

The exact commands, options, and configuration keys are maintained in the
project's build, deployment, and framework documentation. This chapter
describes the workflow and its inputs; it does not specify guarantees for
installation rollback, release promotion, or unattended process recovery.
